\section{An actual Gaussian profile at an even first exponent}
\label{ge-section}

Let $a,b$ be positive coprime integers, and write
\[
 c=a+b,\quad A=ab,\quad M=a^2+ab+b^2=c^2-A,\qquad
 F=a^4+3a^3b+5a^2b^2+3ab^3+b^4.
\]
The elementary primitive identities give
\begin{equation}\label{ge-norm-identities}
 F=M^2+Ac^2=A^2+Mc^2,
 \qquad \gcd(A,c)=\gcd(A,M)=\gcd(M,c)=1,
 \qquad \gcd(F,AMc)=1.
\end{equation}
Also $F$ is odd, prime to $3$, and at least $13$. Every prime divisor
$q$ of $F$ is $1$ modulo $3$: the nonzero residue $M/A$ satisfies
$X^2+X+1=0$ and has exact order $3$, since $q\ne3$.

\begin{theorem}[Quadratic support condition]\label{ge-character}
Suppose $M=DU^2$ with positive integers $D,U$, without requiring $D$ to
be squarefree. Every prime $q\mid F$ satisfies
\begin{equation}\label{ge-quadratic-character}
 q\nmid DUc,\qquad \left(\frac{-D}{q}\right)=1.
\end{equation}
Thus an actual first profile $M=3^eV_0R^h$, where $e\in\{0,1\}$ and
$h\ge2$ is even, imposes on every prime divisor $q$ of the whole second
norm, and hence on every prime divisor of any second root $Q$, the
conditions
\begin{equation}\label{ge-residual-character}
 e=0:\quad\left(\frac{-V_0}{q}\right)=1;
 \qquad
 e=1:\quad\left(\frac{V_0}{q}\right)=1.
\end{equation}
No equality of the two exponents or power-free residual is assumed.
\end{theorem}
\begin{proof}
Equation~\eqref{ge-norm-identities} becomes
$F=A^2+D(Uc)^2$, with $\gcd(A,Uc)=1$. At a prime divisor $q$ of $F$,
the quotient $A/(Uc)$ is nonzero and its square is $-D$.
Set $D=3^eV_0$ and use $(-3/q)=1$, which follows from $q\equiv1\pmod3$,
to obtain~\eqref{ge-residual-character}.
\end{proof}
The corresponding quadratic field has discriminant dividing $4D$;
the logarithm of its absolute value is at most $\log12+\log V_0$.
This gives no uniform prime
distribution estimate while the residual, and therefore the field, moves.

\begin{theorem}[An actual Gaussian decomposition]\label{ge-gaussian-lift}
Assume $M=U^2$ for a positive integer $U$. Then
\begin{equation}\label{ge-gaussian-element}
 z_G=A+iUc\in\mathbb Z[i],\qquad
 \operatorname N(z_G)=F,\qquad \gcd(A,Uc)=1.
\end{equation}
The Gaussian integer is unramified at $2$, and every prime $q\mid F$
satisfies $q\equiv1\pmod{12}$.
For any integers $g\ge2$, $V_1>0$, $Q>1$ with $F=V_1Q^g$, there are
primitive Gaussian integers $v_G,w_G$ and a unit $u_G\in\{1,-1,i,-i\}$
such that
\begin{equation}\label{ge-gaussian-profile}
 z_G=u_Gv_Gw_G^g,\qquad
 \operatorname N(v_G)=V_1,\qquad \operatorname N(w_G)=Q.
\end{equation}
Neither factor contains a pair of conjugate primes or a prime over $2$.
Their supports may overlap; $V_1$ need not be $g$-free.
\end{theorem}
\begin{proof}
Primitivity and the norm identity follow from
\eqref{ge-norm-identities}; $F$ odd excludes the prime over $2$.
The Gaussian Euclidean algorithm follows by choosing a nearest Gaussian
integer to a complex quotient, leaving a remainder of norm at most half
the divisor norm. Thus $\mathbb Z[i]$ is a UFD. An inert rational prime
cannot divide a primitive Gaussian integer, and the two conjugate primes
over a split rational prime cannot both divide it. Hence every $q\mid F$
is $1$ modulo $4$, and also $1$ modulo $3$ as above.

For its uniquely occurring orientation $\pi_q$, the exponent in $z_G$
is exactly
$v_q(F)=v_q(V_1)+g v_q(Q)$. Assign these two nonnegative parts to
$v_G$ and $w_G$. Equality is then up to a Gaussian unit, proving the
decomposition and its stated primitive properties.
\end{proof}
If $V_1=1$, this is the pure power $z_G=u_Gw_G^g$. For odd $g$, the
unit can be absorbed into $w_G$, since the $g$-power map permutes the
four Gaussian units. The unit is retained for a general exponent.

At the same time the actual second Eisenstein element has the reviewed
oriented decomposition
\begin{equation}\label{ge-two-domains}
 z_E=A+U^2\zeta=u_Ev_Ew_E^g,
 \qquad \operatorname N(v_E)=V_1,\quad\operatorname N(w_E)=Q.
\end{equation}
Thus the two powers in different quadratic UFDs share the rational
coordinate $A$, with the further constraint
\begin{equation}\label{ge-shared-coordinate}
 (Uc)^2=U^2(A+U^2).
\end{equation}
Independently chosen powers need not satisfy either this constraint or
equality of their rational coordinates.

\begin{corollary}[Height of the actual Gaussian coefficient]
\label{ge-coefficient-height}
With the data of~\eqref{ge-gaussian-profile}, put
$\eta_G=(u_G^2v_G/\overline v_G)^{-1}$. Then
\[
 \left(\frac{w_G}{\overline w_G}\right)^g
   =\eta_G\frac{z_G}{\overline z_G},
 \qquad h(\eta_G)=\frac12\log V_1,
\]
where $h$ is absolute logarithmic Weil height.
\end{corollary}
\begin{proof}
Conjugating the decomposition proves the identity. There is no
conjugate-prime cancellation in $v_G/\overline v_G$, whose denominator
ideal has norm $V_1$ and whose complex absolute value is one.
The normalized height is therefore $\frac12\log V_1$; roots of unity
and inversion leave it unchanged.
\end{proof}
This is a coefficient-height statement for actual points, with no
point-height upper bound or independence assertion for the two covers.

\begin{theorem}[Cyclotomic equation and exact integer inverse]
\label{ge-cyclotomic-inverse}
If $M=U^2$ for a positive integer $U$, then, with $d=a-b$,
\begin{equation}\label{ge-phi12}
 F=U^4-U^2c^2+c^4=\Phi_{12}(U,c),\qquad
 d^2=4U^2-3c^2,\qquad |d|<c,
\end{equation}
and
\[
 \gcd(U,c)=1,\qquad \frac{\sqrt3}{2}<\frac Uc<1.
\]
Conversely, let $U,c$ be positive coprime integers with $U<c$ and let
$d\in\mathbb Z$ satisfy $d^2=4U^2-3c^2$ and $|d|<c$. Then
\begin{equation}\label{ge-inverse-seed}
 a=(c+d)/2,\qquad b=(c-d)/2
\end{equation}
are positive coprime integers with $a^2+ab+b^2=U^2$, and their second
norm is $\Phi_{12}(U,c)$.
\end{theorem}
\begin{proof}
Use $A=c^2-U^2$ and $d^2=c^2-4A$ in
\eqref{ge-norm-identities}. The upper interval bound follows from $A>0$.
Equality in the lower bound would give $a=b$, hence the primitive seed
$(1,1)$ with nonsquare first norm $3$; it cannot occur here.

For the inverse, $d^2\equiv c^2\pmod4$ gives matching parity, and
$|d|<c$ gives positivity. Substitution proves $a^2+ab+b^2=U^2$ and
$ab=c^2-U^2$. An odd prime dividing both $a,b$ would divide $c,d$ and
therefore $4U^2$, contrary to $\gcd(U,c)=1$. If both $a,b$ are even,
their norm identity forces $U$ even, while $c$ is even, the same
contradiction. This proves coprimality and all the claimed identities.
\end{proof}
The extra square condition in~\eqref{ge-phi12} is part of the inverse.
No identification with arbitrary rational points of
$y^g=\Phi_{12}(x,1)$ is made. To retain a specified larger even first
exponent $h$, one must additionally impose $U=R^{h/2}$ with an integer
$R>1$. The ramified branch $M=3U^2$ instead gives
$F=9U^4-3U^2c^2+c^4$ and has not been replaced by the same integral
cyclotomic equation.

\paragraph{Remaining scope.}
These statements apply through the pure-coefficient inverse theorem to
every finite positive rational point whose reconstructed first norm is
an unramified even power. The even/even exclusion is the earlier
combined-square result. The odd second exponent remains unresolved.
For a pure second prime exponent $p>7$, the earlier modular support
theorem gives, for every prime $q\mid Q$,
$q\ge(\sqrt{2p}-1)^2>p$, now intersected with $q\equiv1\pmod{12}$.
This is a necessary support condition, not control of the large-prime
tail. The ramified first branch and moving nonunit residuals remain
open. The proofs here are ordinary arithmetic proofs, independently
reviewed; no complete geometric or modular Lean formalization is claimed.
